\documentclass[dvisvgm,tikz,border=3pt]{standalone}
\usepackage{amsmath,amssymb}
\usepackage{pgfplots}
\usepackage{CJKutf8}
\pgfplotsset{compat=1.18}
\usetikzlibrary{arrows.meta,calc,positioning}

\definecolor{themebg}{HTML}{2e362c}
\definecolor{themefg}{HTML}{edf4e8}
\definecolor{themegray}{HTML}{a4af9d}
\definecolor{themecurve}{HTML}{7eb6ff}
\definecolor{themealert}{HTML}{ff7b7b}
\definecolor{themeamber}{HTML}{f5b942}

\pgfdeclarelayer{background}
\pgfdeclarelayer{foreground}
\pgfsetlayers{background,main,foreground}

\begin{document}
\begin{CJK*}{UTF8}{gbsn}
\pagecolor{themebg}
\begin{tikzpicture}[color=themefg, text=themefg]

% ── 上图: 被积函数 f(t) 含跳跃间断点 ──
\begin{axis}[
    name=axTop,
    width=11.5cm,
    height=4.8cm,
    axis lines=middle,
    axis line style={color=themefg, -{Stealth}},
    tick style={color=themefg},
    tick label style={color=themefg, font=\scriptsize},
    every axis label/.style={color=themefg},
    xmin=-0.5, xmax=4.5,
    ymin=-0.3, ymax=2.4,
    xtick={1, 2.5},
    xticklabels={$a$, $c$ (跳跃点)},
    ytick={1, 2},
    yticklabels={$h_1$, $h_2$},
    clip=false,
    xlabel={$t$},
    ylabel={$f(t)$},
    title style={font=\footnotesize\bfseries, at={(0.5,1.08)}, anchor=south},
    title={一阶累积的光滑性升阶机制：被积函数跳跃间断 $\implies$ 变限积分连续但出现折角},
]

% 积分阴影区域
\begin{pgfonlayer}{background}
    \fill[themecurve!25!themebg] (axis cs:1,0) rectangle (axis cs:2.5,1);
    \fill[themecurve!35!themebg] (axis cs:2.5,0) rectangle (axis cs:3.5,2);
\end{pgfonlayer}

% 被积函数台阶波形
\addplot[themecurve, line width=2.0pt, domain=1:2.5] {1};
\addplot[themecurve, line width=2.0pt, domain=2.5:4] {2};

% 断点空心与实心标记
\addplot[only marks, mark=*, mark size=2.2pt, themecurve] coordinates {(2.5, 1)};
\addplot[only marks, mark=o, mark size=2.2pt, line width=1.2pt, themecurve, fill=themebg] coordinates {(2.5, 2)};

% 竖直虚线对齐到下方的折角点
\draw[themegray, dashed, line width=1.0pt] (axis cs:2.5, 0) -- (axis cs:2.5, 2.2);

\node[text=themecurve, font=\scriptsize\bfseries, anchor=west] at (axis cs:2.7, 1.0) {$f(t)$ 跳跃间断};
\end{axis}

% ── 下图: 变上限面积函数 F(x) = \int_a^x f(t)dt ──
\begin{axis}[
    name=axBot,
    at={($(axTop.south)-(0, 1.2cm)$)},
    anchor=north,
    width=11.5cm,
    height=4.8cm,
    axis lines=middle,
    axis line style={color=themefg, -{Stealth}},
    tick style={color=themefg},
    tick label style={color=themefg, font=\scriptsize},
    every axis label/.style={color=themefg},
    xmin=-0.5, xmax=4.5,
    ymin=-0.3, ymax=3.6,
    xtick={1, 2.5},
    xticklabels={$a$, $c$ (连续不可导)},
    ytick=\empty,
    clip=false,
    xlabel={$x$},
    ylabel={$F(x)$},
]

% 连续折线 F(x): 在 [1, 2.5] 斜率为 1，在 [2.5, 4] 斜率为 2
\addplot[themeamber, line width=2.2pt, domain=1:2.5] {x - 1};
\addplot[themeamber, line width=2.2pt, domain=2.5:4] {1.5 + 2*(x - 2.5)};

% 竖直对齐虚线
\draw[themegray, dashed, line width=1.0pt] (axis cs:2.5, 0) -- (axis cs:2.5, 1.5);

% 折角尖点标记
\begin{pgfonlayer}{foreground}
    \addplot[only marks, mark=*, mark size=2.8pt, fill=themealert, draw=themefg] coordinates {(2.5, 1.5)};
    \node[text=themealert, font=\scriptsize\bfseries, anchor=south east] at (axis cs:2.45, 1.6) {折角点 $F'_-(c)=h_1 \ne F'_+(c)=h_2$};
    \node[text=themeamber, font=\scriptsize\bfseries, anchor=west] at (axis cs:3.2, 2.2) {$F(x)=\int_a^x f(t)dt$};

    % 光滑性升级结论框
    \node[text=themefg, font=\scriptsize, anchor=north west] at (axis cs:-0.3, 3.4) {
        \begin{tabular}{l}
        \textbf{光滑性升阶规则}：\\
        1. $f(t)$ 可积 $\implies F(x)$ 必\textbf{连续}\\
        2. $f(t)$ 连续 $\implies F(x)$ 必\textbf{可导}且 $F'(x)=f(x)$\\
        3. $f(t)$ $C^k$ 光滑 $\implies F(x)$ 为 $C^{k+1}$ 光滑
        \end{tabular}
    };
\end{pgfonlayer}

\end{axis}

\end{tikzpicture}
\end{CJK*}
\end{document}
